% Original analytic teaching diagram, 2026-11-12.
% Two overlapping events A and B cut Omega into four atoms; sigma({A,B}) consists of the 2^4 = 16 unions of atoms.
\definecolor{VB}{HTML}{7188A3}\definecolor{VBF}{HTML}{E5EAF0}
\definecolor{VG}{HTML}{829C8B}\definecolor{VGF}{HTML}{E8EEE9}
\definecolor{VY}{HTML}{B4A16F}\definecolor{VYF}{HTML}{F1ECDD}
\definecolor{VR}{HTML}{A7787D}\definecolor{VRF}{HTML}{F0E5E6}
\definecolor{VP}{HTML}{9588A5}\definecolor{VPF}{HTML}{ECE8F0}
\definecolor{VI}{HTML}{3B4252}\definecolor{VM}{HTML}{6B7280}\definecolor{VL}{HTML}{4C4D4F}
\begin{tikzpicture}[x=1mm,y=1mm,>=Stealth,font=\figurefont\fontsize{8}{10}\selectfont,
lab/.style={text=VM,font=\figurefont\fontsize{7.5}{9}\selectfont}]
\fill[VPF] (0,0) rectangle (64,32);
\fill[VBF] (24,16) ellipse (16 and 11);
\fill[VYF] (40,16) ellipse (16 and 11);
\begin{scope}
  \clip (24,16) ellipse (16 and 11);
  \fill[VGF] (40,16) ellipse (16 and 11);
\end{scope}
\draw[VB,line width=.8pt] (24,16) ellipse (16 and 11);
\draw[VY,line width=.9pt] (40,16) ellipse (16 and 11);
\draw[VL,line width=.6pt] (0,0) rectangle (64,32);
\node[text=VI] at (16,16) {$A\cap B^{\mathrm{c}}$};
\node[text=VI] at (32,16) {$A\cap B$};
\node[text=VI] at (48,16) {$A^{\mathrm{c}}\cap B$};
\node[text=VI] at (54,3.5) {$A^{\mathrm{c}}\cap B^{\mathrm{c}}$};
\node[text=VB] at (13,25.5) {$A$};
\node[text=VY] at (51,25.5) {$B$};
\node[text=VI] at (3.5,29) {$\Omega$};
\end{tikzpicture}
